\addtocontents{toc}{\protect\Needspace{3\baselineskip}}
\originalchapter{第二次测验复习：曲率、体积与闭宇宙光路}
8.286 Fall 2013 第二次测验复习册的几何问题多次变换径向坐标。已有广义相对论卷提供度规、测地线的前置知识；这里针对原题的可测长度、光行积分和参数分支给出完整题解，不重新铺陈广义相对论体系。

\src{8.286 Quiz 2 Review Problem 4, 原题为1990年 Quiz 2；官方复习解答。}
\begin{exercise}
开放、尘埃主导模型的参数解为 $a/\sqrt\kappa=\alpha(\cosh\theta-1)$、$ct=\alpha(\sinh\theta-\theta)$，$\alpha>0$。求 $a/\sqrt\kappa=2\alpha$ 时的宇宙时刻。
\end{exercise}
\begin{solution}
从 $\cosh\theta-1=2$ 得膨胀分支 $\theta=\operatorname{arcosh}3$，$\sinh\theta=\sqrt{3^2-1}=2\sqrt2$，故
\[
t=\frac{\alpha}{c}\left(2\sqrt2-\operatorname{arcosh}3\right)
\simeq1.06567\,\frac\alpha c .
\]
取负 $\theta$ 会对应相反时间分支，不符合大爆炸后膨胀的题设。
\end{solution}

\src{8.286 Quiz 2 Review Problem 5；官方复习解答。}
\begin{exercise}
若宇宙是闭合、仅尘埃的模型，今天测得 $\Omega_0=2$、$H_0>0$，离大坍缩还剩多少时间？用 $H_0$ 表示，明确选择膨胀分支。
\end{exercise}
\begin{solution}
闭宇宙参数解可写 $a/\sqrt{k}=\alpha(1-\cos\theta)$，$ct=\alpha(\theta-\sin\theta)$，其中 $0<\theta<\pi$ 是膨胀支，终点 $\theta=2\pi$。由
$\Omega_0-1=kc^2/(a_0^2H_0^2)=1$ 得 $a_0/\sqrt{k}=c/H_0$。代入尘埃参数常数
$\alpha=4\pi G\rho_0a_0^3/(3k^{3/2}c^2)$，及 $\rho_0=2(3H_0^2/8\pi G)$，得 $\alpha=c/H_0$。于是 $1-\cos\theta_0=1$ 且取膨胀支 $\theta_0=\pi/2$；现在年龄 $(\pi/2-1)/H_0$，总寿命 $2\pi/H_0$，剩余
\[
t_{\rm crunch}-t_0=\frac{1+3\pi/2}{H_0}.
\]
若误选收缩支 $3\pi/2$，会把今天仍有正 $H_0$ 的条件丢掉。
\end{solution}

\src{8.286 Quiz 2 Review Problem 6(a)--(d), 原题为1998年 Quiz 2 Q3；官方复习解答。}
\begin{exercise}
闭合 FLRW 模型在径向角坐标 $\psi$ 下有
$\dd s^2=-c^2\dd t^2+a(t)^2[\dd\psi^2+\sin^2\psi\,\dd\Omega^2]$。
(a) 求径向光的 $\dd\psi/\dd t$；(b) 用积分表示时刻 $t$ 的粒子视界径向行程实际长度；(c) 对尘埃参数解 $ct=\alpha(\theta-\sin\theta)$、$a=\alpha(1-\cos\theta)$ 求 $\dd\psi/\dd\theta$；(d) 一束光在大爆炸极限从 $\psi=0$ 出发，绕闭空间一周返回原点相当于宇宙完整寿命的多少？
\end{exercise}
\begin{solution}
零线元给 $\dd\psi/\dd t=\pm c/a(t)$；从初始极限到 $t$ 的\emph{沿光路}实际径向行程为
$a(t)\int_0^t c\,\dd t'/a(t')$。在尘埃参数式中 $\dd t/\dd\theta=a/c$，所以向外光线 $\dd\psi/\dd\theta=1$。闭三维球面上的一条大圆需角行程 $2\pi$ 才回原点；参数 $\theta$ 从 0 到 $2\pi$ 也恰是大爆炸至大坍缩，所求寿命比为 1，按奇点端点取极限理解。
当累计光路角超过 $\pi$，它已越过对跖点；公式中的 $a\psi$ 是沿原定方向的\emph{行程}，不总是该时刻两点间最短测地距离。这一点是使用“视界实际长度”时的坐标边界。
\end{solution}

\src{8.286 Quiz 2 Review Problem 7(a)--(b), 原题为1994年 Quiz 2 Q3；官方复习解答。}
\begin{exercise}
二维空间度规 $\dd s^2=(1+\alpha r)^2\dd r^2+r^2(1+\beta r)^2\dd\theta^2$，$\alpha,\beta>0$。从原点沿 $\theta=0$ 到 $r_0$，再沿 $r=r_0$ 到 $\theta=\pi/2$，最后沿固定角回原点。(a) 求路径总长；(b) 求所围实际面积。
\end{exercise}
\begin{solution}
两条径向边各长 $\int_0^{r_0}(1+\alpha r)\dd r=r_0+\alpha r_0^2/2$；弧长为 $(\pi/2)r_0(1+\beta r_0)$。总长
$S=(2+\pi/2)r_0+(\alpha+\pi\beta/2)r_0^2$。
面积元是 $\sqrt{\det g}\,\dd r\dd\theta=r(1+\alpha r)(1+\beta r)\dd r\dd\theta$，因此
\[
A=\frac{\pi}{2}\left[\frac{r_0^2}{2}
+\frac{(\alpha+\beta)r_0^3}{3}
+\frac{\alpha\beta r_0^4}{4}\right].
\]
不能把坐标扇形的欧氏面积 $\pi r_0^2/4$ 当作弯曲度规下的全部答案。
\end{solution}

\src{8.286 Quiz 2 Review Problem 8(a)--(c), 原题为1988年 Quiz 2 Q4；官方复习解答。}
\begin{exercise}
在某固定宇宙时刻，闭合 FLRW 空间度规为
$\dd\ell^2=a^2[\dd r^2/(1-r^2)+r^2(\dd\theta^2+\sin^2\theta\,\dd\phi^2)]$。
一小杆的两端在 $(h,0,0)$ 与 $(h,\Delta\theta,0)$，$\Delta\theta\ll1$。(a) 求原点到第一端的实际径向距离 $\ell_p$；(b) 求杆长 $s_p$；(c) 消去 $h$，用 $\ell_p,s_p,a$ 表示张角。
\end{exercise}
\begin{solution}
在此坐标支 $0\le h\le1$，
$\ell_p=a\int_0^h\dd r/\sqrt{1-r^2}=a\arcsin h$，
$s_p=ah\Delta\theta$。因 $h=\sin(\ell_p/a)$，
\[
\Delta\theta=\frac{s_p}{a\sin(\ell_p/a)}.
\]
这是小杆极限；$r=\sin\psi$ 在越过赤道后不单值，若用另一支须改用完整 $\psi$ 坐标并处理透镜/多光路，不把 $\arcsin$ 主值延拓到整颗三维球面。
\end{solution}

\src{8.286 Quiz 2 Review Problem 9(a)--(d), 原题为1986年 Quiz 2 Q3；官方复习解答。}
\begin{exercise}
三维球对称空间采用实际径向坐标 $r$，度规
$\dd\ell^2=\dd r^2+\varrho(r)^2(\dd\theta^2+\sin^2\theta\,\dd\phi^2)$。
对 $r\le r_0$：(a) 求实际半径；(b) 求边界球面积；(c) 求包含体积；(d) 若换坐标 $\sigma=r^2$，写新度规。
\end{exercise}
\begin{solution}
$r$ 的径向线元本来就是 $\dd r$，故半径为 $r_0$；固定 $r_0$ 面的面积是 $4\pi\varrho(r_0)^2$。体积元 $\sqrt{\det g}=\varrho(r)^2\sin\theta$，故
$V=4\pi\int_0^{r_0}\varrho(r)^2\dd r$。在 $r=\sqrt\sigma$、$\sigma>0$ 的坐标支，
\[
\dd\ell^2=\frac{\dd\sigma^2}{4\sigma}
+\varrho(\sqrt\sigma)^2
(\dd\theta^2+\sin^2\theta\,\dd\phi^2).
\]
$\sigma=0$ 的系数发散是新坐标在原点退化，不代表原几何必有曲率发散。
\end{solution}

\src{8.286 Quiz 2 Review Problem 10, 原题为1990年 Quiz 3 Q1；官方复习解答。}
\begin{exercise}
某固定时刻的 Robertson--Walker 空间度规
$\dd\ell^2=a^2[\dd r^2/(1-kr^2)+r^2(\dd\theta^2+\sin^2\theta\,\dd\phi^2)]$。
对原点中心坐标球 $0\le r\le r_{\max}$ 求实际体积，角积分需做完，径向积分可保留。
\end{exercise}
\begin{solution}
行列式平方根为 $a^3r^2\sin\theta/\sqrt{1-kr^2}$，所以
\[
V(r_{\max})=4\pi a^3\int_0^{r_{\max}}
\frac{r^2\dd r}{\sqrt{1-kr^2}}.
\]
积分区间须在所用坐标支内满足 $1-kr^2>0$。闭空间 $k>0$ 时此 $r$ 坐标对完整三维球面并非全局单值，不能把到 $r=1/\sqrt k$ 的单支积分当作全空间体积；改用 $\psi$ 支才可计两半。
\end{solution}
